\subsection{The Model Context Protocol}
MCP standardizes how LLM agents discover and invoke external tools. A client and
server negotiate a protocol version and capabilities during an \texttt{initialize}
handshake. The client can then list tools (\texttt{tools/list}) and invoke them
(\texttt{tools/call}). Messages follow JSON-RPC 2.0, including its error-object
semantics and its reserved error codes. Most deployments run the server locally over
stdio, where stdout carries protocol traffic and nothing else. The specification is
versioned by date, and servers may negotiate down to a version they support.

\subsection{Static studies of the MCP ecosystem}
The ecosystem has been studied at scale, but rarely by running it.
Health-and-maintainability work analyzes source and metadata for roughly 1,900
servers using static scanners and code-smell detectors~\cite{mcpfirstglance}, while
registry measurement crawls several marketplaces and characterizes scale,
maintenance and supply-chain concentration across some 8,000 servers from marketplace data, without running them~\cite{mcpcrawler}. A companion dataset effort catalogues 3,238 candidate MCP
repositories on GitHub and classifies them by operational role, again from repository evidence~\cite{mcpgithubdataset}, and a parallel line studies
the \emph{application} side, characterizing how 1,723 MCP applications configure and
supervise the servers they consume~\cite{mcpapps}. Fault taxonomies take a third
route and mine GitHub issues~\cite{runtimefaults,realfaults}. One finds tool faults
and data-schema faults the two most frequent categories~\cite{runtimefaults}, the same two areas our harness tests directly. None of this
work executes a server.

\subsection{Execution-based work is security-focused}
Some prior work does execute MCP servers, to find security problems. Sandboxed behavioral scanners hunt vulnerabilities in curated
corpora~\cite{sandboxscan}, taint-style analyzers confirm exploitable flaws across
tens of thousands of repositories~\cite{viper}, a security-invariants benchmark avoids live
third-party execution altogether in favor of fixtures and a 40-repository metadata
sample~\cite{securityinvariants}, and malicious-server detection, authentication
studies, and dynamic assessment of internet-facing endpoints scope to security
surfaces~\cite{maliciousmcp,remoteauth,exposedbydesign}.

MCPZoo assembles 64,611 unique servers, of
which more than 37,288 support dynamic analysis, and uses them to show that existing
security scanners are unreliable, with fewer than half of sampled alerts holding up
under manual validation~\cite{mcpzoo}. That corpus differs from ours in two ways. First, MCPZoo is built by \emph{repairing}
servers: a multi-agent framework infers environments and iteratively fixes deployment
and runtime defects so that in-the-wild repositories become running services. That
construction is the right choice for studying what a server does once it runs, and it
necessarily removes the signal we measure, which is how often a server runs at all
when a client installs it as published. Second, its population is collected from
multiple public markets, whereas ours is the official registry, the surface a client actually
resolves against. That their corpus needed a repair framework fits our runnability finding.
Conformance is a separate question: whether a server that is neither malicious nor vulnerable does what the
specification says.

Concurrently with our work, Afsar draws a seeded random sample of 400 npm stdio
servers from the same registry and probes each without repair: 48.8\% complete a
handshake, servers that never start (37.5\%) outnumber those missing credentials
(13.3\%), and the 195 that run show no fatal schema violations across 2,766
tools~\cite{randomdraw2026}. That study samples one package ecosystem and checks
declared structure. We census both npm and PyPI and test behavior, calling tools
with inputs the server should reject.

\subsection{The specification's own error-handling debate}
The categorization this paper measures is under active revision. SEP-1303, accepted
and reflected in the current specification, moved tool input validation errors out of
protocol errors and into tool execution errors, reasoning that a language model can
self-correct only from an error that reaches its context~\cite{sep1303}. SEP-2145,
open at the time of writing, extends the same treatment to tool resolution failures,
including unknown tools, and states its purpose as aligning the specification with
existing Python and TypeScript SDK behavior~\cite{sep2145}. We take no position on which mechanism is correct. We measure what
fraction of the deployed ecosystem the proposed alignment already describes.

Concurrent with this work, a practitioner census posted to the same proposal probed
roughly 10,500 \emph{remote} registry endpoints, harvested their tool lists, and
reported that \(88.8\%\) of mutating tools declare no output schema at all, arguing
that the proposal's output-validation clause is inert for most of the write
surface~\cite{siliroid2026census}. That measurement and ours are close to disjoint.
It probes hosted HTTP endpoints and reads declared metadata, whereas we install and
execute local stdio servers and test behavior against inputs. It addresses the
output-validation clause, and we address tool resolution and input enforcement.

\subsection{Conformance tooling and the gap we fill}
The protocol maintainers ship a conformance suite for testing client and server
implementations, built for SDK maintainers to run in continuous integration; to our
knowledge it has not been used to measure the deployed ecosystem. As far as we know, ours is the first reproducible conformance measurement of this ecosystem
that installs and executes the servers themselves and tests their behavior against
adversarial inputs, released with an open harness and dataset.
